\begin{figure}[ht]
    \centering
    \begin{tikzpicture}[
        >=latex,
        block/.style={
            rectangle,
            draw,
            minimum width=1cm,
            minimum height=1cm
        },
        sum/.style={
            circle,
            draw,
            minimum size=0.7cm
        },
        integral/.style={
            rectangle,
            draw,
            minimum width=1cm,
            minimum height=1cm
        }
    ]
    % Input and junction point
    \node (input) at (-6,0) {$u$};
    \node[circle,inner sep=1pt,fill=black] (junction1) at (-5,0) {};
    
    % Main forward path
    \node[block] (B) at (-3.5,0) {$B$};
    \node[sum] (sum1) at (-2,0) {$+$};
    \node[integral] (integral) at (-0.5,0) {$\int$};
    \node[block] (C) at (2,0) {$C$};
    \node[sum] (sum2) at (3.5,0) {$+$};
    \node (output) at (5,0) {$y$};
    
    % Feedback path with A
    \node[block] (A) at (0,-2) {$A$};
    
    % Feed-forward path with D
    \node[block] (D) at (0,2) {$D$};
    
    % State labels
    \node at (-1.2,0.4) {$\dot{x}$};
    \node at (0.5,0.4) {$x$};
    
    % Connections
    \draw[->] (input) -- (junction1);
    \draw[->] (junction1) -- (B);
    \draw[->] (B) -- (sum1);
    \draw[->] (sum1) -- (integral);
    \draw[->] (integral) -- (C);
    \draw[->] (C) -- (sum2);
    \draw[->] (sum2) -- (output);
    
    % Feedback connection
    \draw[->] (integral) -- (integral |- A.north) -- (A);
    \draw[->] (A) -- (A -| sum1.south) -- (sum1.south);
    
    % Feed-forward connection
    \draw[->] (junction1) |- (D);
    \draw[->] (D) -| (sum2);
    
    \end{tikzpicture}
    \caption{State Space Model Diagram}
    \label{fig:ssm}
\end{figure}